% ---------------------------------------------------------------------
% Parte 9 · Lecciones y cierre · expone César · 0:30 (lo edita César)
% Tres lecciones concretas para sus proyectos, una por tarjeta, y una
% frase de cierre. Nada genérico: cada lección debe salir de lo expuesto.
% ---------------------------------------------------------------------
\begin{frame}{Lo que nos llevamos}
  \begin{columns}[T]
    \begin{column}{0.31\textwidth}
      \begin{tarjeta}[equal height group=lecciones]{bcmTCP}{1}
        \small Usar \textbf{TCP} cuando los datos no pueden perderse ni llegar en desorden, como en APIs, bases de datos y transacciones.
      \end{tarjeta}
    \end{column}
    \begin{column}{0.31\textwidth}
      \begin{tarjeta}[equal height group=lecciones]{bcmUDP}{2}
        \small Elegir \textbf{UDP} si conviene evitar el handshake y las retransmisiones, asumiendo que la app tolera pérdidas o las maneja aparte.
      \end{tarjeta}
    \end{column}
    \begin{column}{0.31\textwidth}
      \begin{tarjeta}[equal height group=lecciones]{bcmWS}{3}
        \small Aplicar \textbf{WebSockets} en paneles en vivo o chats donde el servidor debe enviar datos continuos sin abrir nuevas peticiones.
      \end{tarjeta}
    \end{column}
  \end{columns}
  \vfill
  \begin{center}
    {\large\bfseries En un sistema distribuido real los protocolos no compiten: se combinan según lo que necesita cada módulo.\par}
  \end{center}
  \vfill
\end{frame}